
% --- PACOTES BÁSICOS E IDIOMA ---
\usepackage[ruled, linesnumbered, portuguese]{algorithm2e}
\usepackage{amsmath}
\usepackage[utf8]{inputenc}
\usepackage[T1]{fontenc}
\usepackage[portuguese]{babel} % Garante nomes como "Sumário" e hifenização correta
\usepackage{graphicx}
\usepackage{xcolor} % Para controle de cores nos l
\usepackage{float} % no preâmbulo

% --- MATEMÁTICA AVANÇADA ---
\usepackage{amsmath, amssymb, amsthm, bm}
\usepackage{esint} % Para integrais avançadas (como a da sua imagem inicial)

% --- FORMATAÇÃO DE PÁGINA (ABNT) ---
\usepackage[top=3cm, bottom=2cm, left=3cm, right=2cm]{geometry} % Margens padrão
\usepackage{setspace}
\onehalfspacing % Espaçamento de 1.5
\usepackage{indentfirst} % Indenta o primeiro parágrafo de cada seção

% --- TABELAS E FORMATAÇÃO ---
\usepackage{tabularx}
\usepackage{booktabs}
\usepackage{array}

% --- ORGANIZAÇÃO E REFERÊNCIAS ---
\usepackage{makeidx}
\usepackage[acronym]{glossaries}
\usepackage{hyperref}
\hypersetup{
	colorlinks=true,
	linkcolor=black,   % Links internos em preto para impressão
	citecolor=black,   % Citações em preto
	urlcolor=blue,     % Links da web em azul
	pdftitle={Caracterização Óptica de Meios Coloidais}
}

% --- DEFINIÇÕES DE TEOREMAS ---
\theoremstyle{remark}
\newtheorem{Remark}{Remark}

\theoremstyle{definition}
\newtheorem{Example}{Exemplo}
\newtheorem{Theorem}{Teorema}
\newtheorem{Definition}{Definição}
\newtheorem{Lemma}{Lema}
\newtheorem{corollary}{Corolário}